% ==============================================================================
% PREAMBOLO — Stile accademico LaTeX classico (sobrio, bianco e nero standard)
% ==============================================================================
\usepackage[T1]{fontenc}
\usepackage[utf8]{inputenc}
\usepackage[italian]{babel}
\usepackage{csquotes}
\usepackage{lmodern}
\usepackage{microtype}

% Geometria pagina standard bilanciata
\usepackage[
  a4paper,
  top=2.5cm,
  bottom=2.5cm,
  left=2.4cm,
  right=2.4cm
]{geometry}
\usepackage{placeins}

% Matematica e unita' di misura (standard italiano con virgola decimale)
\usepackage{amsmath, amssymb, mathtools}
\usepackage{siunitx}
\sisetup{
  locale = DE,
  output-decimal-marker = {,},
  group-separator = {\,},
  group-minimum-digits = 4,
  per-mode = symbol,
  range-phrase = {--},
  range-units = single,
  detect-all = true
}
\DeclareSIUnit\bit{bit}
\DeclareSIUnit\byte{B}
\DeclareSIUnit\spike{spike}
\DeclareSIUnit\event{evento}
\DeclareSIUnit\events{eventi}
\DeclareSIUnit\ch{ch}
\DeclareSIUnit\hop{hop}
\DeclareSIUnit\class{class.}

% Palette neutrale in scala di grigi (stile classico LaTeX da stampa)
\usepackage[dvipsnames, svgnames, table]{xcolor}
\definecolor{NavyPrimary}{gray}{0.10}
\definecolor{TealAccent}{gray}{0.25}
\definecolor{CoralAccent}{gray}{0.20}
\definecolor{AmberAccent}{gray}{0.35}
\definecolor{SlateText}{gray}{0.15}
\definecolor{LightBg}{gray}{0.96}
\definecolor{WarmBg}{gray}{0.93}
\definecolor{TealBg}{gray}{0.94}
\definecolor{RuleGray}{gray}{0.75}

% Sfondi neutri molto tenui per i blocchi circuitali
\definecolor{BlockNMDA}{gray}{0.94}
\definecolor{BlockLEAK}{gray}{0.91}
\definecolor{BlockAHP}{gray}{0.95}
\definecolor{BlockNa}{gray}{0.92}
\definecolor{BlockK}{gray}{0.96}
\definecolor{BlockRed}{gray}{0.93}

% Grafica, figure, sottotitoli e tabelle standard
\usepackage{graphicx}
\graphicspath{{figures/}}
\usepackage{subcaption}
\usepackage{float}
\usepackage{booktabs}
\usepackage{tabularx}
\newcolumntype{L}[1]{>{\raggedright\arraybackslash}p{#1}}
\newcolumntype{Y}{>{\raggedright\arraybackslash}X}
\usepackage{multirow}
\usepackage{makecell}
\usepackage{pdflscape}
\usepackage{rotating}
\usepackage{enumitem}
\setlist{nosep, leftmargin=*, topsep=4pt}

% Didascalie standard LaTeX in nero
\usepackage[
  font=small,
  labelfont=bf,
  margin=6pt
]{caption}

% Circuiti CMOS (CircuiTikZ), diagrammi a blocchi (TikZ) e grafici (PGFPlots)
\usepackage[siunitx, RPvoltages]{circuitikz}
\ctikzset{
  transistors/scale=0.82,
  resistors/scale=0.78,
  capacitors/scale=0.78,
  sources/scale=0.78,
  amplifiers/scale=0.78,
  logic ports/scale=0.75,
  tripoles/mos style/arrows
}
\usetikzlibrary{
  arrows.meta,
  positioning,
  calc,
  fit,
  backgrounds,
  shapes.geometric,
  shapes.multipart,
  decorations.pathreplacing
}
\usepackage{pgfplots}
\pgfplotsset{compat=1.18}
\usepgfplotslibrary{groupplots}

% Riquadri di sintesi sobri in stile classico bianco/grigio
\usepackage[most]{tcolorbox}
\newtcolorbox{sintesiBox}[1]{
  enhanced,
  breakable,
  colback=gray!4,
  colframe=black!70,
  colbacktitle=gray!12,
  coltitle=black,
  fonttitle=\bfseries\small,
  title={#1},
  boxrule=0.5pt,
  titlerule=0.4pt,
  arc=1.5pt,
  left=8pt, right=8pt, top=6pt, bottom=6pt,
  before skip=10pt, after skip=10pt
}

\newtcolorbox{criticaBox}[2][]{
  enhanced,
  breakable,
  colback=gray!4,
  colframe=black!80,
  colbacktitle=gray!15,
  coltitle=black,
  fonttitle=\bfseries\small,
  title={#2},
  boxrule=0.6pt,
  titlerule=0.4pt,
  arc=1.5pt,
  left=8pt, right=8pt, top=6pt, bottom=6pt,
  before skip=10pt, after skip=10pt,
  #1
}

\newtcolorbox{notaTecnica}[1]{
  enhanced,
  breakable,
  colback=white,
  colframe=black!60,
  colbacktitle=gray!10,
  coltitle=black,
  fonttitle=\bfseries\small,
  title={#1},
  boxrule=0.5pt,
  titlerule=0.4pt,
  arc=1.5pt,
  left=8pt, right=8pt, top=6pt, bottom=6pt,
  before skip=8pt, after skip=8pt
}

% Stile di pagina standard LaTeX (numero di pagina pulito in basso, zero decorazioni)
\pagestyle{plain}

% Bibliografia con BibLaTeX + Biber (stile numerico standard)
\usepackage[
  backend=biber,
  style=numeric-comp,
  sorting=none,
  maxbibnames=6,
  minbibnames=3,
  giveninits=true,
  doi=true,
  url=false,
  isbn=false
]{biblatex}
\addbibresource{bibliografia.bib}

% Collegamenti ipertestuali cliccabili ma completamente neri (stile standard)
\usepackage[
  hidelinks,
  bookmarksopen=true,
  bookmarksnumbered=true
]{hyperref}

% Scorciatoie tipografiche utili e coerenti
\newcommand{\Vmem}{V_{\mathrm{mem}}}
\newcommand{\Vthr}{V_{\mathrm{thr}}}
\newcommand{\Vrst}{V_{\mathrm{rst}}}
\newcommand{\Vrest}{V_{\mathrm{rest}}}
\newcommand{\Vdd}{V_{\mathrm{DD}}}
\newcommand{\Cmem}{C_{\mathrm{mem}}}
\newcommand{\Csyn}{C_{\mathrm{syn}}}
\newcommand{\Isyn}{I_{\mathrm{syn}}}
\newcommand{\Imem}{I_{\mathrm{mem}}}
\newcommand{\Iin}{I_{\mathrm{in}}}
\newcommand{\UT}{U_{\mathrm{T}}}
\newcommand{\ND}{\text{N.D.}}

% Spaziatura paragrafi pulita e leggibile
\setlength{\parindent}{0pt}
\setlength{\parskip}{5pt}
\renewcommand{\arraystretch}{1.18}
